% Options for packages loaded elsewhere
\PassOptionsToPackage{unicode}{hyperref}
\PassOptionsToPackage{hyphens}{url}
\PassOptionsToPackage{dvipsnames,svgnames,x11names}{xcolor}
\documentclass[
  10pt,
  fontset=fandol]{ctexart}
\usepackage{xcolor}
\usepackage[margin=16mm]{geometry}
\usepackage{amsmath,amssymb}
\setcounter{secnumdepth}{-\maxdimen} % remove section numbering
\usepackage{iftex}
\ifPDFTeX
  \usepackage[T1]{fontenc}
  \usepackage[utf8]{inputenc}
  \usepackage{textcomp} % provide euro and other symbols
\else % if luatex or xetex
  \usepackage{unicode-math} % this also loads fontspec
  \defaultfontfeatures{Scale=MatchLowercase}
  \defaultfontfeatures[\rmfamily]{Ligatures=TeX,Scale=1}
\fi
\usepackage{lmodern}
\ifPDFTeX\else
  % xetex/luatex font selection
\fi
% Use upquote if available, for straight quotes in verbatim environments
\IfFileExists{upquote.sty}{\usepackage{upquote}}{}
\IfFileExists{microtype.sty}{% use microtype if available
  \usepackage[]{microtype}
  \UseMicrotypeSet[protrusion]{basicmath} % disable protrusion for tt fonts
}{}
\makeatletter
\@ifundefined{KOMAClassName}{% if non-KOMA class
  \IfFileExists{parskip.sty}{%
    \usepackage{parskip}
  }{% else
    \setlength{\parindent}{0pt}
    \setlength{\parskip}{6pt plus 2pt minus 1pt}}
}{% if KOMA class
  \KOMAoptions{parskip=half}}
\makeatother
\usepackage{longtable,booktabs,array}
\newcounter{none} % for unnumbered tables
\usepackage{calc} % for calculating minipage widths
% Correct order of tables after \paragraph or \subparagraph
\usepackage{etoolbox}
\makeatletter
\patchcmd\longtable{\par}{\if@noskipsec\mbox{}\fi\par}{}{}
\makeatother
% Allow footnotes in longtable head/foot
\IfFileExists{footnotehyper.sty}{\usepackage{footnotehyper}}{\usepackage{footnote}}
\makesavenoteenv{longtable}
\setlength{\emergencystretch}{3em} % prevent overfull lines
\providecommand{\tightlist}{%
  \setlength{\itemsep}{0pt}\setlength{\parskip}{0pt}}
\usepackage{amsmath,amssymb,mathtools,needspace}
\xeCJKDeclareCharClass{CJK}{"2032,"0400->"04FF,"2160->"216B,"2460->"2473,"25A0,"25C7,"25CB,"25CF}
\IfFontExistsTF{Arial Unicode MS}{\xeCJKsetup{AutoFallBack=true}\setCJKfallbackfamilyfont{\CJKrmdefault}{Arial Unicode MS}}{}
\setcounter{secnumdepth}{0}
\setlength{\emergencystretch}{2em}
\allowdisplaybreaks
\usepackage{bookmark}
\IfFileExists{xurl.sty}{\usepackage{xurl}}{} % add URL line breaks if available
\urlstyle{same}
\hypersetup{
  pdftitle={奇光异彩},
  colorlinks=true,
  linkcolor={Maroon},
  filecolor={Maroon},
  citecolor={Blue},
  urlcolor={Blue},
  pdfcreator={LaTeX via pandoc}}

\title{奇光异彩}
\author{}
\date{}

\begin{document}
\maketitle

\subsection{12.3　奇光异彩}\label{ux5947ux5149ux5f02ux5f69}

我们看到，尽管二分割圆生成的多边形面积 \(S_n\) 逼近于圆面积 \(S^*\)，但逼近过程收

\href{../../source-images/pdf-298.jpeg}{原页图像}

敛缓慢，为了获得高精度的圆周率，所要耗费的计算量可能变得很大。譬如，需要割到 24 576 边形，才能得出祖冲之的``密率''3.141 592 6。在古代用算筹一类简单的计算工具，实现如此浩大的计算工程是难以想象的。

\subsubsection{12.3.1　刘徽的新视野}\label{ux5218ux5fbdux7684ux65b0ux89c6ux91ce}

面对这个现实，刘徽独具慧眼地提出了这样一个挑战性的课题：设法将已经获得的数据进行``再加工''，希望以尽量少的计算量为代价获得高精度的结果。

刘徽用一个具体的算例

\[
\widehat S=S_{192}+\frac{36}{105}(S_{192}-S_{96})\approx S_{3\,072}
\]

演示了这种设计方案的可行性。

被称为``刘徽神算''的这个数学案例，实际上表达了一种数据精加工的方法。推广刘徽神算，自然可提出下述\textbf{校正技术}：

设法寻求某个\textbf{松弛因子} \(\omega\)，将偏差 \(\Delta_n=S_{2n}-S_n\) 的 \(\omega\) 倍作为数据 \(S_{2n}\) 的\textbf{校正量}，而使\textbf{校正值}

\[
\widehat S=S_{2n}+\omega(S_{2n}-S_n),\quad 0<\omega<1
\tag{12.3.1}
\]

比 \(S_n\) 和 \(S_{2n}\) 具有更高的精度。

自然会问，作为校正技术的式（12.3.1），为什么要选取形如 \(\omega(S_{2n}-S_n)\) 的校正项呢？

前已指出刘徽在《割圆术》中导出了圆面积的双侧挤压公式

\[
S_{2n}<S^*<S_{2n}+(S_{2n}-S_n),
\]

借助于偏差 \(\Delta_n=S_{2n}-S_n\) 这个公式可表达为

\[
0\cdot\Delta_n<S^*-S_{2n}<1\cdot\Delta_n,
\]

即误差 \(S^*-S_{2n}\) 被夹在系数分别为 0 和 1 的左右两极之间。

中华文化崇尚``中庸之道''，无论处理什么事情都要避免走极端，尽量不左不右，不偏不倚，执中致和。因此，依据上述挤压公式，应令

\[
S^*-S_{2n}\approx\omega\Delta_n,
\]

式中 \(0<\omega<1\)，即取式（12.3.1）作为校正公式。

另一方面，校正公式（12.3.1）亦可改写成两个数据 \(S_n\) 与 \(S_{2n}\) 的组合形式：

\[
\widehat S=(1+\omega)S_{2n}-\omega S_n.
\]

按这种理解，所述精加工方法也是一种\textbf{组合技术}。

问题在于，这里组合系数为什么采取一正一负的逆反形式呢？

事实上，比较 \(S_n\) 和 \(S_{2n}\) 两个近似值，虽然它们都很粗糙，但 \(S_{2n}\) 总比 \(S_n\) 更为精确，即 \(S_{2n}\) 为``优''值而 \(S_n\) 则为``劣''值，所以加权平均应采取``激浊扬清''的态势，充分激发 \(S_{2n}\) 的``优''势而抑制 \(S_n\) 的``劣''势，为此令 \(S_n\) 的权系数 \(-\omega<0\)，而令 \(S_{2n}\) 的权系数 \(1+\omega>1\)。

\href{../../source-images/pdf-299.jpeg}{原页图像}

激浊扬清、优劣互补的态势，标志着这种技术从属于高效算法设计的二分演化模式。

总之，\textbf{作为刘徽神算推广的校正公式（12.3.1），是校正技术与组合技术两种技术的综合。}为使这类技术真正保证提高精度的要求，该怎样具体地选取松弛因子 \(\omega\) 呢？

\subsubsection{12.3.2　偏差比中传出好``消息''}\label{ux504fux5deeux6bd4ux4e2dux4f20ux51faux597dux6d88ux606f}

再换一种视角考察前述数据校正技术。针对无穷逼近过程 \(\{S_n\}\)，能否设计更``好''的逼近过程 \(\{\widehat S_n\}\)，使之具有更快的收敛速度呢？这就是逼近加速问题。

欲使\textbf{校正公式}（12.3.1）成为逼近加速公式，校正因子 \(\omega\) 应当是驾驭逼近过程的某个数学不变量，即与逼近过程相关的某个普适常数，称之为\textbf{加速因子}。前已指出，这个常数 \(\omega\) 应当在 0 与 1 之间，即 \(0<\omega<1\)。

数学常数通常采取比率的形式，刘徽指出，比率的本意是相关事物的比例关系。

在二分割圆过程中，需要考虑什么样的``相关量''呢？

《割圆术》称偏差 \(\Delta_n=S_{2n}-S_n\) 为\textbf{差幂}，刘徽特别关注\textbf{幂率} \(\delta_n=\Delta_n/\Delta_{2n}\)。在割圆计算过程中，利用差幂 \(\Delta_n\) 计算幂率 \(\delta_n\)，割圆术实际上已造出下列数据表 12.1。

\textbf{表 12.1}

\begingroup
\renewcommand{\arraystretch}{2.2}

{\def\LTcaptype{none} % do not increment counter
\begin{longtable}[]{@{}rrrr@{}}
\toprule\noalign{}
\(n\) & \(S_n\) & \(\Delta_n\) & \(\delta_n\) \\
\midrule\noalign{}
\endhead
\bottomrule\noalign{}
\endlastfoot
12 & 300 & \(10\dfrac{364}{625}\) & 3.95 \\
24 & \(310\dfrac{364}{625}\) & \(2\dfrac{425}{625}\) & 3.99 \\
48 & \(313\dfrac{164}{625}\) & \(\dfrac{420}{625}\) & 4.00 \\
96 & \(313\dfrac{584}{625}\) & \(\dfrac{105}{625}\) & \\
192 & \(314\dfrac{64}{625}\) & & \\
\end{longtable}
}

\endgroup

从这些数据中能获得什么样的信息呢？无需用高深的数学知识进行理论分析，\textbf{直接观察幂率 \(\delta_n\) 的数据表（见表 12.1）即可发现，幂率几乎为定值 4。}

\textbf{刘徽由此发现了一个奇妙的事实：表面上杂乱无章的数据 \(S_n\) 中竟蕴藏着极其鲜明的规律性。}

\subsubsection{12.3.3　只要做一次``俯冲''}\label{ux53eaux8981ux505aux4e00ux6b21ux4fefux51b2}

偏差比几乎是个定值，基于这个奇妙的规律，分析误差只是一蹴而就的事，据此只要做一次``俯冲''便能捕捉到所要的校正公式。

事实上，由于在割圆计算过程中偏差比近似等于 4，从而得出一系列近似关系

\href{../../source-images/pdf-300.jpeg}{原页图像}

式：

\[
\begin{aligned}
S_{2n}-S_n&\approx4(S_{4n}-S_{2n}),\\
S_{4n}-S_{2n}&\approx4(S_{8n}-S_{4n}),\\
S_{8n}-S_{4n}&\approx4(S_{16n}-S_{8n}),\\
&\vdots\\
S_{2N}-S_N&\approx4(S_{4N}-S_{2N}),
\end{aligned}
\]

式中，\(N\) 是某个远大于 \(n\) 的正整数。

将这些式子累加在一起，其中间项相互抵消，得

\[
S_{2N}-S_n\approx4(S_{4N}-S_{2n}).
\]

这样，若取 \(S_{4N}\) 和 \(S_{2N}\) 作为 \(S_{2n}\) 的校正值 \(\widehat S\)，则有

\[
\widehat S-S_n=4(\widehat S-S_{2n}),
\tag{12.3.2}
\]

即有

\[
\frac{\widehat S-S_n}{\widehat S-S_{2n}}\approx4.
\]

如果称近似值 \(S_n\) 与校正值 \(\widehat S\) 两者之差为\textbf{残差}，则上式说明，\textbf{若偏差比几乎为定值 4，那么残差比也近似等于 4。}

这样一来，精加工方法的设计就水到渠成了。事实上，从式（12.3.2）解出未知的 \(\widehat S\)，有

\[
\widehat S=S_{2n}+\frac{1}{3}(S_{2n}-S_n).
\tag{12.3.3}
\]

这个式子的含义是，在偏差比``几乎''为定值 4 的情况下，近似值的残差比也近似等于 4，因而残差 \(\widehat S-S_{2n}\) 几乎等于偏差 \(S_{2n}-S_n\) 的 \(\dfrac13\)，从而立即导出所求的加速公式（12.3.3）。

值得指出的是，刘徽神算（12.2.2）即

\[
\widehat S=S_{192}+\frac{36}{105}(S_{192}-S_{96})
\]

中的校正因子 \(\omega=\dfrac{36}{105}\) 非常接近实际精加工公式（12.3.3）中的校正因子 \(\omega=\dfrac13=\dfrac{35}{105}\)。至于刘徽为什么这样处理，\textbf{一个明显的理由是，刘徽过于偏爱计算公式与计算结果的简洁美。}

\subsubsection{12.3.4　差之毫厘，失之千里}\label{ux5deeux4e4bux6bebux5398ux5931ux4e4bux5343ux91cc}

加速因子究竟怎样选取才算合适呢？循着刘徽割圆术的思路，精确地计算直

\href{../../source-images/pdf-301.jpeg}{原页图像}

到正 3 072 边形的面积，直接选取加速因子 \(\omega=\dfrac13\)，再按式（12.3.3）计算，计算结果列于表 12.2 中，表中括弧 \(\langle\,\cdot\,\rangle\) 标明数据准确到小数点后第几位。\(\pi\) 的真值为 3.141 592 653 5\ldots。

\textbf{表 12.2}

{\def\LTcaptype{none} % do not increment counter
\begin{longtable}[]{@{}
  >{\raggedleft\arraybackslash}p{(\linewidth - 4\tabcolsep) * \real{0.3333}}
  >{\raggedleft\arraybackslash}p{(\linewidth - 4\tabcolsep) * \real{0.3333}}
  >{\raggedleft\arraybackslash}p{(\linewidth - 4\tabcolsep) * \real{0.3333}}@{}}
\toprule\noalign{}
\begin{minipage}[b]{\linewidth}\raggedleft
\(n\)
\end{minipage} & \begin{minipage}[b]{\linewidth}\raggedleft
\(S_n\)
\end{minipage} & \begin{minipage}[b]{\linewidth}\raggedleft
\(\widehat S_n\)
\end{minipage} \\
\midrule\noalign{}
\endhead
\bottomrule\noalign{}
\endlastfoot
12 & \(3.000\,000\,000\quad\langle0\rangle\) & \\
24 & \(3.105\,828\,541\quad\langle1\rangle\) & \(3.141\,104\,722\quad\langle3\rangle\) \\
48 & \(3.132\,628\,613\quad\langle1\rangle\) & \(3.141\,561\,971\quad\langle4\rangle\) \\
96 & \(3.139\,350\,203\quad\langle1\rangle\) & \(3.141\,590\,733\quad\langle5\rangle\) \\
192 & \(3.141\,031\,951\quad\langle3\rangle\) & \(3.141\,592\,534\quad\langle6\rangle\) \\
384 & \(3.141\,452\,472\quad\langle3\rangle\) & \(3.141\,592\,646\quad\langle7\rangle\) \\
768 & \(3.141\,557\,608\quad\langle4\rangle\) & \(3.141\,592\,653\quad\langle8\rangle\) \\
1 536 & \(3.141\,583\,892\quad\langle4\rangle\) & \(3.141\,592\,654\quad\langle8\rangle\) \\
3 072 & \(3.141\,590\,463\quad\langle5\rangle\) & \(3.141\,592\,654\quad\langle8\rangle\) \\
\end{longtable}
}

在割圆计算中，刘徽已获知直到正 3 072 边形的数据。以上数据表显示，利用这些数据按照刘徽加速技术进行精加工，即可获得祖冲之的``密率''3.141 592 65。

刘徽如果这样做，那么中华数学史乃至世界数学史上有关圆周率科学计算部分就要彻底改写了。

\textbf{差之毫厘，失之千里。刘徽的校正技术极为精彩，只是由于校正因子的处理稍欠精细，从而失之交臂，把一项千年称雄的数学成就留给了两百年后的祖冲之。}

\subsubsection{12.3.5　``缀术''再剖析}\label{ux7f00ux672fux518dux5256ux6790}

本章一开头指出，南北朝数学家祖冲之给出了高精度的圆周率，这项成就领先世界一千多年。

祖冲之的圆周率是怎样得出来的？如果直接用内接正多边形来逼近，要一直算到正 \(24\,576=6\times2^{12}\) 边形。耗费如此巨大的计算量，在筹算的古代是难以想象的。

中华民族是个智慧的民族，是个善于创新的民族。\textbf{刘徽的加速技术是中华文明前瞻性思维的一个明证。}

据隋唐古书记载，祖冲之的算法设计技术称为``缀术''。史书盛赞缀术``指要精密，

\begin{quote}
校注（agent 补充）：表 12.2 的 \(S_n\) 数值从 3.000 000 000 开始，按原表保留；此前半径 \(r=10\) 寸的表 12.1 数值从 300 开始。本页未明说此处改用单位圆或作了归一化，转录未擅自改变尺度。
\end{quote}

\href{../../source-images/pdf-302.jpeg}{原页图像}

算氏之最''。但关于缀术的具体内容已无从考证，仅仅靠``缀术''这个名词能够窥探出其中的奥秘吗？

查看汉语词典，\textbf{``缀''字有两个含义：一是``缀合''，即组合，因此``缀术''就是组合技术；另一个含义是``缀补''，即修补和校正，在这个意义上，``缀术''亦可理解为校正技术。}

\textbf{这样，所谓``缀术''其实就是组合技术与校正技术。}本节一开头就指出，刘徽的精加工技术正是这种技术。也许，据此可以断定，\textbf{祖冲之的``缀术''正是继承了刘徽的``衣钵''------特别是刘徽的割圆术，归纳总结出来的。}

数学史先辈钱宝琮先生早就猜测过祖冲之与刘徽之间的传承关系。我们深信，\textbf{祖冲之将自己的算法设计技术命名为``缀术''，目的也是试图向世人表白：自己的数学成就，只是前人特别是刘徽的研究工作的缀补和修正。因此，``缀术''实质上只是《九章算术》刘徽注的祖冲之注。}

也许这就是历史的真相。

\subsubsection{12.3.6　平庸的新纪录}\label{ux5e73ux5eb8ux7684ux65b0ux7eaaux5f55}

前已指出，祖冲之于公元 5 世纪所获得的准确到小数点后 7 位的圆周率 3.141 592 6\ldots 千年称雄于世界。这项纪录直到 15 世纪才被打破。阿拉伯人阿尔·卡西于 1424 年写成《圆周论》，发表了当时世界上最为精确的圆周率。

阿尔·卡西的割圆过程袭用阿基米德的做法：从正六边形做起，逐步计算圆的内接与外切多边形的周长。每分割一次，令多边形的边数倍增。到了 15 世纪，阿拉伯数字和十进小数记数法的使用，给实际计算提供了很大方便。阿尔·卡西反复割圆，一直算到

\[
6\times2^{27}=805\,306\,368
\]

即正 8 亿多多边形，得出准确到小数点后 17 位的圆周率

\[
\pi=3.141\,592\,653\,589\,793\,23,
\]

从而打破了祖冲之千年称雄的世界纪录。

阿尔·卡西的这项计算其实是平庸的。其一，他所使用的其实就是刘徽公式；其二，后文将会看到，运用刘徽的加速技术，用正 24 576 边形的数据就能加工出阿尔·卡西正 8 亿多多边形的结果。

这里再现阿尔·卡西所获得的数据。令直径为 1，则内接正 \(n\) 边形的边长

\[
L_n=n\cdot\sin\frac{\pi}{n}.
\]

设取 \(n=6,12,24,\cdots\) 反复进行计算，计算结果列于表 12.3 中。表中数据是借助于数学软件 Mathematica 获得的，计算过程避开了舍入误差的积累。

\begin{quote}
校注（agent 补充）：本页圆周率末段印作 \(\ldots793\,23\)，PDF294（书页281）同一阿尔·卡西结果印作 \(\ldots792\,23\)，两处原文不一致，均按原页保留。

校注（agent 补充）：原文称 \(L_n=n\sin(\pi/n)\) 为``边长''，该式实际为直径 1 的圆内接正 \(n\) 边形的周长；单边长应为 \(\sin(\pi/n)\)。这是原书文字与所列公式的不一致，正文未作替换。
\end{quote}

\href{../../source-images/pdf-303.jpeg}{原页图像}

\textbf{表 12.3}

{\def\LTcaptype{none} % do not increment counter
\begin{longtable}[]{@{}
  >{\raggedleft\arraybackslash}p{(\linewidth - 4\tabcolsep) * \real{0.3333}}
  >{\raggedleft\arraybackslash}p{(\linewidth - 4\tabcolsep) * \real{0.3333}}
  >{\raggedleft\arraybackslash}p{(\linewidth - 4\tabcolsep) * \real{0.3333}}@{}}
\toprule\noalign{}
\begin{minipage}[b]{\linewidth}\raggedleft
\(n\)
\end{minipage} & \begin{minipage}[b]{\linewidth}\raggedleft
\(L_n\)
\end{minipage} & \begin{minipage}[b]{\linewidth}\raggedleft
\(\delta_n\)
\end{minipage} \\
\midrule\noalign{}
\endhead
\bottomrule\noalign{}
\endlastfoot
6 & \(3.000\,000\,000\,000\,000\,000\quad\langle0\rangle\) & \(3.948\,815\,549\,036\,689\,1\) \\
12 & \(3.105\,828\,541\,230\,249\,148\quad\langle1\rangle\) & \(3.987\,162\,707\,851\,017\,4\) \\
24 & \(3.132\,628\,613\,281\,238\,197\quad\langle1\rangle\) & \(3.996\,788\,098\,191\,334\,7\) \\
48 & \(3.139\,350\,203\,046\,867\,207\quad\langle1\rangle\) & \(3.999\,196\,863\,295\,489\,2\) \\
96 & \(3.141\,031\,950\,890\,509\,638\quad\langle3\rangle\) & \(3.999\,799\,205\,744\,363\,9\) \\
192 & \(3.141\,452\,472\,285\,462\,075\quad\langle3\rangle\) & \(3.999\,949\,800\,806\,102\,4\) \\
384 & \(3.141\,557\,607\,911\,857\,645\quad\langle4\rangle\) & \(3.999\,987\,450\,162\,151\,0\) \\
768 & \(3.141\,583\,892\,148\,318\,408\quad\langle4\rangle\) & \(3.999\,996\,862\,538\,076\,8\) \\
1 536 & \(3.141\,590\,463\,228\,050\,095\quad\langle5\rangle\) & \(3.999\,999\,215\,634\,365\,4\) \\
3 072 & \(3.141\,592\,105\,999\,271\,550\quad\langle6\rangle\) & \(3.999\,999\,803\,908\,581\,7\) \\
6 144 & \(3.141\,592\,516\,692\,157\,447\quad\langle6\rangle\) & \(3.999\,999\,950\,977\,144\,8\) \\
12 288 & \(3.141\,592\,619\,365\,383\,955\quad\langle7\rangle\) & \(3.999\,999\,987\,744\,286\,1\) \\
24 576 & \(3.141\,592\,645\,033\,690\,896\quad\langle7\rangle\) & \(3.999\,999\,996\,936\,071\,5\) \\
49 152 & \(3.141\,592\,651\,450\,767\,651\quad\langle8\rangle\) & \(3.999\,999\,999\,234\,017\,8\) \\
98 304 & \(3.141\,592\,653\,055\,036\,841\quad\langle9\rangle\) & \(3.999\,999\,999\,808\,504\,4\) \\
196 608 & \(3.141\,592\,653\,456\,104\,139\quad\langle9\rangle\) & \(3.999\,999\,999\,952\,126\,1\) \\
393 216 & \(3.141\,592\,653\,556\,370\,963\quad\langle10\rangle\) & \(3.999\,999\,999\,988\,031\,5\) \\
786 432 & \(3.141\,592\,653\,581\,437\,669\quad\langle11\rangle\) & \(3.999\,999\,999\,997\,007\,8\) \\
1 572 864 & \(3.141\,592\,653\,587\,704\,346\quad\langle11\rangle\) & \(3.999\,999\,999\,999\,251\,9\) \\
3 145 728 & \(3.141\,592\,653\,589\,271\,015\quad\langle12\rangle\) & \(3.999\,999\,999\,999\,812\,9\) \\
6 291 456 & \(3.141\,592\,653\,589\,662\,682\quad\langle12\rangle\) & \(3.999\,999\,999\,999\,953\,2\) \\
12 582 912 & \(3.141\,592\,653\,589\,760\,599\quad\langle13\rangle\) & \(3.999\,999\,999\,999\,988\,3\) \\
25 165 824 & \(3.141\,592\,653\,589\,785\,078\quad\langle13\rangle\) & \(3.999\,999\,999\,999\,997\,0\) \\
50 331 648 & \(3.141\,592\,653\,589\,791\,198\quad\langle14\rangle\) & \(3.999\,999\,999\,999\,999\,2\) \\
100 663 296 & \(3.141\,592\,653\,589\,792\,728\quad\langle14\rangle\) & \(3.999\,999\,999\,999\,999\,8\) \\
201 326 592 & \(3.141\,592\,653\,589\,793\,110\quad\langle15\rangle\) & \(3.999\,999\,999\,999\,999\,9\) \\
402 653 184 & \(3.141\,592\,653\,589\,793\,206\quad\langle16\rangle\) & \\
805 306 368 & \(3.141\,592\,653\,589\,793\,230\quad\langle17\rangle\) & \\
\end{longtable}
}

\href{../../source-images/pdf-304.jpeg}{原页图像}

表 12.3 所列的计算结果表明，在二分逼近过程中，反复计算内接正 \(n\) 边形的周长 \(L_n\)，发现阿尔·卡西的``新纪录''果然是对的。

为了运用刘徽方法进行精加工，首先需要澄清偏差比是否``几乎''为定值。利用逼近数据 \(L_n\) 计算偏差 \(\Delta_n=L_{2n}-L_n\)，进而求出偏差比 \(\delta_n=\Delta_n/\Delta_{2n}\)，计算结果 \(\delta_n\) 列于表 12.3 的右侧。我们看到，这里偏差比 \(\delta_n\) 越来越逼近定值 4，因此加速公式仍具有式（12.3.3）的形式：

\[
\widehat L_n=L_{2n}+\frac13(L_{2n}-L_n).
\]

依这一公式加工表 12.3 的数据 \(L_n\)，加工结果 \(\widehat L_n\) 列于表 12.4 中。数据尾部依然标明准确到小数第几位。

\textbf{表 12.4}

{\def\LTcaptype{none} % do not increment counter
\begin{longtable}[]{@{}
  >{\raggedleft\arraybackslash}p{(\linewidth - 6\tabcolsep) * \real{0.0700}}
  >{\raggedleft\arraybackslash}p{(\linewidth - 6\tabcolsep) * \real{0.4300}}
  >{\raggedleft\arraybackslash}p{(\linewidth - 6\tabcolsep) * \real{0.0700}}
  >{\raggedleft\arraybackslash}p{(\linewidth - 6\tabcolsep) * \real{0.4300}}@{}}
\toprule\noalign{}
\begin{minipage}[b]{\linewidth}\raggedleft
\(n\)
\end{minipage} & \begin{minipage}[b]{\linewidth}\raggedleft
\(\widehat L_n\)
\end{minipage} & \begin{minipage}[b]{\linewidth}\raggedleft
\(n\)
\end{minipage} & \begin{minipage}[b]{\linewidth}\raggedleft
\(\widehat L_n\)
\end{minipage} \\
\midrule\noalign{}
\endhead
\bottomrule\noalign{}
\endlastfoot
6 & \(3.141\,104\,721\,640\,332\,197\quad\langle3\rangle\) & 768 & \(3.141\,592\,653\,587\,960\,658\quad\langle11\rangle\) \\
12 & \(3.141\,561\,970\,631\,567\,880\quad\langle4\rangle\) & 1 536 & \(3.141\,592\,653\,589\,678\,702\quad\langle12\rangle\) \\
24 & \(3.141\,590\,732\,968\,743\,543\quad\langle5\rangle\) & 3 072 & \(3.141\,592\,653\,589\,786\,079\quad\langle13\rangle\) \\
48 & \(3.141\,592\,533\,505\,057\,115\quad\langle6\rangle\) & 6 144 & \(3.141\,592\,653\,589\,792\,791\quad\langle14\rangle\) \\
96 & \(3.141\,592\,646\,083\,779\,554\quad\langle7\rangle\) & 12 288 & \(3.141\,592\,653\,589\,793\,210\quad\langle16\rangle\) \\
192 & \(3.141\,592\,653\,120\,656\,168\quad\langle9\rangle\) & 24 576 & \(3.141\,592\,653\,589\,793\,236\quad\langle17\rangle\) \\
384 & \(3.141\,592\,653\,560\,471\,996\quad\langle10\rangle\) & & \\
\end{longtable}
}

对照表 12.3 和表 12.4 可以看出，\textbf{为了获得阿尔·卡西割到正 8 亿多多边形所创造的新纪录，运用刘徽的精加工方法只要割到正 24 576 边形，后者相当于祖冲之已经掌握的逼近数据。}

\textbf{刘徽加速技术的效果简直令人难以置信！}

\begin{center}\rule{0.5\linewidth}{0.5pt}\end{center}

\subsection{相关原页校注（agent 补充）}\label{ux76f8ux5173ux539fux9875ux6821ux6ce8agent-ux8865ux5145}

以下校注来自本节覆盖的原页，可能也涉及同页相邻小节；原文照录与校正说明分开保留。

\subsubsection{原书 PDF 页 304 的校注}\label{ux539fux4e66-pdf-ux9875-304-ux7684ux6821ux6ce8}

\href{../pages/pdf-304.md}{完整原页转录} · \href{../../source-images/pdf-304.jpeg}{原页图像}

校注记录：CH12-E004。

\begin{quote}
校注（agent 补充）：按本页原式，表中 \(\widehat L_{24\,576}\) 需要 \(L_{24\,576}\) 和 \(L_{49\,152}\) 两个数据。正文``只要割到正 24 576 边形''与公式及表下标所表示的数据需求不一致；原文和数值均予保留。
\end{quote}

\subsection{本节覆盖原页的校注记录}\label{ux672cux8282ux8986ux76d6ux539fux9875ux7684ux6821ux6ce8ux8bb0ux5f55}

下列记录按原页关联，含同页相邻小节的校注；计算时同时读取上文校注和完整原页。

\begin{itemize}
\tightlist
\item
  \texttt{CH12-E002}：原书 PDF 页 294, 302
\item
  \texttt{CH12-E003}：原书 PDF 页 302
\item
  \texttt{CH12-E004}：原书 PDF 页 302, 304
\end{itemize}

\end{document}
